Une onde transversale sur une corde est décrite par
\[ y(x,t) = \num{\v{A}}\cos\left(\num{\v{k}}\,x - \num{\v{w}}\,t\right) \]
où $y$ est en centimètres, $x$ en mètres, $t$ en secondes et la phase en radians. On s'intéresse au point de la corde situé à $x = \qty{\v{x1}}{m}$, à l'instant $t = \qty{\v{t1}}{s}$.

\begin{enumerate}[label=\alph*)]
    \item Déterminez le déplacement de ce point.
    \item Déterminez sa vitesse transversale. Le point monte-t-il ou descend-il ?
    \item Déterminez son accélération transversale.
\end{enumerate}
